An incoming graph has its upper endpoint at infinite radius. We now
put a maximum beyond the inner polynomial and return those graphs to a
finite upper section. A common launch must pair the physical turning
heights $h$ and $h+s+o(1)$, with the same $s$ at every prescribed test.
Write $d$ for the inner degree and $R$ for the crest radius scale. The
scaled height is the physical height divided by $R^d$, and
$\tilde\eps=R^{2-2d}$. On the regular descent the slope of the
action with respect to scaled height is bounded away from zero.
Thus the usual $O(\sqrt{\tilde\eps})$ error for one leg allows a
physical height error of order $R^d\sqrt{\tilde\eps}=R$. To resolve a
fixed shift, Lemma~\ref{lem:deficit-ratio} subtracts the logarithms of
two deficits measured from the same actual reference leg. The common
part of the action then cancels.

Lu matches upper hits on the descent and obtains the first-order
correction corresponding to $U(a)-f(a)\Tr(t)/f(t)$
\cite[supplement, p.~16, (S5.12)--(S5.13)]{LuDegree}.
Two accuracies are needed here. Before the reference hits agree, the
exponential rate locates each upper deficit. After matching, the two
terms in the derivative of the lower-intercept difference have a
common exponential rate through the crest. Their first-order log
ratio determines the sign away from the crest.

\begin{lemma}[Single-half deficits and a common exponent across a crest]\label{lem:common-crest}
Use the scaled legs of Lemma~\ref{lem:variational}, with potential
coordinate $u=r^2/2$ and half-profiles $p_\sigma(r)=\phi_\sigma(u)$,
$\sigma=\pm1$. Let $t_{b,\sigma}\to t_b$ be later reference radii and
fix a finite $T_{\max}>t_b^2/2$. On $\mathcal I=[0,T_{\max}]$ set
$\widetilde E(u)=E(\sqrt{2u})$ and $\widetilde Q(u)=Q(\sqrt{2u})$.
Assume the profiles vanish at the axis, $E(0)=Q(0)=0$, and, as $\delta\downarrow0$,
\[
 \|\phi_\sigma-\widetilde E-\sigma\delta\widetilde Q\|_{C^1(\mathcal I)}
     =o(\delta),\qquad
 \sup_{\delta,\sigma}\operatorname{Var}_{\mathcal I}(\phi_\sigma')<\infty,
 \qquad \sqrt\eps+\eps|\log\eps|=o(\delta).
\]
On each compact subinterval of $(0,T_{\max}]$ assume
$\phi_\sigma\to\widetilde E$ in $C^2$.
For $0<r\le\sqrt{2T_{\max}}$, the even $E$ rises to a simple maximum
at $1$ and then decreases through zero at $1<g<t_b$, with nonzero
slope off the crest. The perturbed profiles likewise have a regular
ascent and descent outside an initial interval $0\le u\le\eta_\delta$,
where $\eta_\delta\to0$ and
$\int_0^{\eta_\delta}|\phi_\sigma'|^2\dd u=o(\delta)$.
Assume transverse axis hits for the turns used below, including the turns
up to each reference endpoint $T_{b,\sigma}=t_{b,\sigma}^2/2$.
Put $f=E'/r$, $U(t)=\int_0^t(fQ'-f'Q)\dd r$, and $\Tr(t)=U(t)-U(1)$.
\begin{enumerate}
  \item[\textup{(i)}] \emph{Single-half deficit rate.} Without assuming
  equality of the reference upper hits, write
  $\mathcal A_{\mathrm{out},\sigma}(S)=\int_0^S \phi_\sigma'\dd m_{\sigma,S}$.
  Uniformly on compact subsets of $T\in[0,t_b^2/2)$ bounded away from $T_{b,\sigma}$,
  \[
    -\eps\log\bigl(X_\sigma(T_{b,\sigma})-X_\sigma(T)\bigr)
    = \min_{S\in[T,T_{b,\sigma}]} \mathcal A_{\mathrm{out},\sigma}(S)
      + O\bigl(\sqrt\eps+\eps|\log\eps|\bigr).
  \]
  In particular, for a descending turn $T$ bounded away from the perturbed
  crest $r_{\mathrm{crest},\sigma}$ and from $T_{b,\sigma}$, the minimum is
  attained at $S=T$ and equals
  $\int_{r_{\mathrm{crest},\sigma}}^{\sqrt{2T}}(\partial_r p_\sigma)^2\dd r/r$.
  \item[\textup{(ii)}] \emph{Matched reference upper hits.} Assume in addition
  that $X_+(T_{b,+})=X_-(T_{b,-})$. Let $t_-(t)$ match the upper hit of the
  plus turn $t$, write $\Delta_Y(t)=Y_+(t^2/2)-Y_-(t_-^2/2)$ and
  $J_\pm=J_{i,\pm}+J_{o,\pm}$ at these actual turns, and for $0<t<g$ let $a=t$
  on the ascent and let $a\le 1$ satisfy $E(a)=E(t)$ on the descent. On compact
  subintervals of $(0,g)$, including those containing $1$, both positive terms
  \[
   \mathcal L=e^{-J_{i,+}},\qquad \mathcal N=e^{J_{o,+}-J_{o,-}-J_{i,-}},\qquad
   \Delta_Y'=\eps t(\mathcal L-\mathcal N)
  \]
  satisfy $-\eps\log \mathcal L,\,-\eps\log \mathcal N\to W(t)=\int_0^a f(r)^2r\dd r$
  uniformly. On compact sets of regular turns the finer limit is
  \begin{equation}\label{eq:crest-splitting}
   \frac{\eps}{2\delta}\log(\mathcal N/\mathcal L)\longrightarrow
   \begin{cases}U(t),&0<t<1,\\
   U(a)-f(a)\Tr(t)/f(t),&1<t<g .
   \end{cases}
  \end{equation}
  If in addition $\Delta_Y(a_*)=0$ for $0<a_*<1$, the signs of the ordinary
  displacement at $a_i<a_*$ and at its descending partner
  $z_i>1$ with $E(z_i)=E(a_i)$ are the negative and the positive,
  respectively, of the corresponding nonzero expressions in
  \eqref{eq:crest-splitting}. A nonzero $U(a_*)$ makes the marked orbit simple.
\end{enumerate}
\end{lemma}

\begin{proof}
For part~\textup{(i)}, suppress $\sigma$ and write $T_b=T_{b,\sigma}$.
For each $0<S\le T_b$ the two envelope actions are
\[
 \mathcal A_{\rm out}(S)=\int_0^S\phi'\dd m_S,\qquad
 I(S)=\int_0^S\phi'\dd l_S,
\]
where $m_S,l_S$ are the suffix extrema in \eqref{eq:running-extrema}.
Integration by parts on $[0,S]$, using \eqref{eq:turn-derivs} and
\eqref{eq:envelope-bounds}, bounds the errors in $-\eps J_o(S)$ and
$\eps J_i(S)$ by $\sqrt{2\eps S}\,(2\|\phi'\|_{\infty,\mathcal I}+
\operatorname{Var}_{\mathcal I}\phi')$. Thus one bound applies from the axis to
either reference turn for all sufficiently small $\delta$.
The envelopes change by at most $\|\phi'\|_\infty |T-T'|$ on their
common interval. The actions are uniformly Lipschitz in the turn by
the same integration by parts. They converge uniformly to those of $E$:
the envelopes converge in sup norm, their variations are bounded, and
the limiting slope is of bounded variation. The limiting upper action
is zero up to the crest and then equals $\int_1^t f^2 r\dd r$, strictly
increasing. The lower action for $0<t<g$ is $W(t)$.

Put $X_*=X_\sigma(T_b)$.
The exact deficit $X_*-X_\sigma(T)=\int_T^{T_b}\eps e^{J_{o,\sigma}(S)}\dd S$
has negative scaled logarithm within
$O(\sqrt\eps+\eps|\log\eps|)$ of $\min_{[T,T_b]}\mathcal A_{\rm out}(S)$,
uniformly away from $T_b$. For the upper bound on the integral use the minimum cost on the
whole interval. For the lower bound integrate on a segment of length
$\eps$ adjacent to a minimizer, where Lipschitz continuity bounds the
increase in action by $O(\eps)$. At a regular descending turn the
suffix maximum follows the descent from the perturbed crest, so the
minimum outgoing cost occurs at the first endpoint and equals
$\int_{r_{\mathrm{crest},\sigma}}^{\sqrt{2T}}(\partial_r p_\sigma)^2\dd r/r$.
This proves part~\textup{(i)}.

For part~\textup{(ii)}, differentiate the equality of the upper hits
$X_+(t^2/2)=X_-(t_-(t)^2/2)$ using \eqref{eq:turn-derivs} to obtain the
formula for $\Delta_Y'$. Part~\textup{(i)} distinguishes the branches for
paired upper hits. On the ascent the limiting upper
height increases to $E(1)$, which rules out a partner beyond the crest
for a regular rising turn. On the descent the limiting deficit cost is
strictly increasing and positive, whereas its value is zero on the ascent.
For a compact test range in $(0,g)$ choose a fixed descending radius
between $g$ and $t_b$. Its limiting deficit cost exceeds all the test
costs, so monotonicity of $X_-$ keeps the partner below this radius.
As $t\to1$, a limiting partner must attain both upper height $E(1)$
and zero deficit cost. A radius below $1$ misses the former and a
radius above $1$ misses the latter. Hence $t_-(t)-t\to0$ uniformly
through the crest as well. The costs of $\mathcal L,\mathcal N$ are therefore
$I_+$ and $I_-+\mathcal A_{{\rm out},+}-\mathcal A_{{\rm out},-}$ up to $O(\sqrt\eps)$, both with limit $W$.

Write $\kappa_+=p_+(t)$ and $\kappa_-=p_-(t_-(t))$ for the
actual scaled turning ordinates of the two arcs with common upper hit.
For a fixed level $\kappa$ on the rising branch, the endpoint defined
by $E(a)+\delta Q(a)=\kappa$ has derivative $-Q(a)/E'(a)$. In the variation of
$\int (E'+\delta Q')^2\dd r/r$ its boundary term cancels $f(a)Q(a)$
in $2\int fQ'$. This gives $U(a)$ up to the initial endpoint term,
which vanishes here ($Q(0)=0$ for the half-profiles).
The analogous descending integral from the crest gives $\Tr(t)$, since
$f(1)=0$. On the initial interval a suffix envelope follows the profile or is
constant, so its action there is bounded by
$\int_0^{\eta_\delta}|\phi_\sigma'|^2\dd u=o(\delta)$.
The $C^1(\mathcal I)$ expansion gives the same bound for the contribution of
$(\widetilde E')^2$ there. Local $C^2$ convergence at the simple crest
controls the perturbed maximum, while the first variations use the
$C^1(\mathcal I)$ remainder.
On the ascent the matched terminal levels $\kappa_\pm$ agree to
$O(\sqrt\eps)=o(\delta)$, so the inward variation gives the first
line of \eqref{eq:crest-splitting}.

On a compact interval of regular descending levels write
$O_\sigma(k)$ for the outgoing envelope action at the turn of level $k$.
Write $z(k)>1$ for the descending inverse of $E$. At zero perturbation
$O_0(k)=\int_1^{z(k)} f(r)^2r\dd r$, hence $O_0'(k)=f(z(k))<0$.
The fixed-level variation just computed gives
$O_\sigma(k)=O_0(k)+\sigma\delta\Tr(z(k))+o(\delta)$ uniformly there.
At the paired turns $T_+=t^2/2$, $T_-=t_-(t)^2/2$, both the upper
hit and the later reference hit agree. Hence the two actual deficits
are the same positive number
\[
 \mathfrak d_*=X_+(T_{b,+})-X_+(T_+)=X_-(T_{b,-})-X_-(T_-).
\]
On a compact set of regular descending turns the pairs stay a fixed
distance from their respective reference endpoints. The outgoing
envelope action on $[T_\sigma,T_{b,\sigma}]$ is minimized at
$T_\sigma$. With $e_\eps=\sqrt\eps+\eps|\log\eps|$, the single-half
deficit estimate above thus gives
$O_\sigma(\kappa_\sigma)=-\eps\log\mathfrak d_*+O(e_\eps)$ for each sign,
uniformly on this set. The bound for the actual outgoing leg now gives
\[
 -\eps\bigl(J_{o,+}(T_+)-J_{o,-}(T_-)\bigr)
 = O_+(\kappa_+)-O_-(\kappa_-)+O(\sqrt\eps)
 = O(e_\eps)=o(\delta).
\]
Since $|O_0'|$ is bounded away from zero, the mean value theorem first
gives $\kappa_+-\kappa_-=O(\delta)$. Both levels tend to $E(t)$, so
applying the same expansion at the two moving levels gives
\[
 o(\delta)=O_+(\kappa_+)-O_-(\kappa_-)
 =f(t)(\kappa_+-\kappa_-)+2\delta\Tr(t)+o(\delta).
\]
Thus $\eps\log(\mathcal N/\mathcal L)=\eps(J_{i,+}-J_{i,-})+o(\delta)$.
The inward action retains the change of level through its derivative
$f(a)$, in addition to its fixed-level variation $U(a)$. Consequently
\[
 \frac{\eps}{2\delta}\log\frac{\mathcal N}{\mathcal L}
 =U(a)+f(a)\frac{\kappa_+-\kappa_-}{2\delta}+o(1)
 =U(a)-f(a)\frac{\Tr(t)}{f(t)}+o(1).
\]
All remainders in this first-order calculation are uniform on compact
sets of regular turns.

On $[a_i,a_*]$ the cost $W$ increases strictly. On $[a_*,z_i]$
it first increases to $W(1)$ and then decreases strictly to
$W(z_i)=W(a_i)<W(a_*)$, since $E(z_i)=E(a_i)<E(a_*)$.
Each test is therefore the unique minimum on its interval to the mark. In a fixed neighborhood of that endpoint the larger of
$\mathcal L,\mathcal N$ exceeds the smaller by $\exp(c\delta/\eps)$ for some $c>0$.
Outside that neighborhood the limiting cost is at least
$W(\text{test})+3\gamma$ for some $\gamma>0$. On a shorter interval
at the test it is below $W(\text{test})+\gamma$.
Uniform logarithmic convergence makes the latter signed mass dominate.
Integrating in the appropriate direction and using
$\sgn\oscD=-\sgn \Delta_Y$ proves the assertions. At the mark the multiplier
is $e^{J_--J_+}$.
\end{proof}

Use the deep orbit of each half as a reference for all its earlier
turns. Lu compares secants to a common terminal section by bounding
the variation of its derivative \cite[Lemma~3.4]{LuDegree}.
On a regular descent we can instead integrate the exact difference of
the two graph equations against the same reference leg. The remainder
in the unscaled logarithm is bounded uniformly in the two turns.

\begin{lemma}[Upper deficits on a regular descent]\label{lem:deficit-ratio}
Let $\eps\downarrow0$. For each scaled profile let $T_b$ be a later turn
and $I=[T_1,T_2]\subset(0,T_b)$ a compact interval of earlier
turns. All these outgoing legs have transverse positive upper hits.
Write $\mathfrak d(T)=X(T_b)-X(T)$. Suppose $|\phi'|\le M$ on $[0,T_b]$,
and on $[T_1,T_b]$ suppose $\nu=-\phi'$ satisfies
$0<\nu_*\le\nu\le\nu^*$ and $|\nu'|\le M_1$. Assume
$T_b-T_2\ge b>0$, $|I|\le L$, and assume the upper hits of the
reference and test legs lie in a compact subinterval of $(0,\infty)$.
For a family of profiles take all these bounds common to the family.
Then, uniformly for $T,T_0\in I$ and all sufficiently small $\eps$,
\begin{equation}\label{eq:deficit-ratio}
 \log\frac{\mathfrak d(T)}{\mathfrak d(T_0)}=-\frac{1}{\eps}\int_{T_0}^{T}
 (\phi')^2\dd u+O(1).
\end{equation}
\end{lemma}
\begin{proof}
Let $w_*$ be the outgoing leg of $T_b$, $\zeta=w_*-\phi$ and
$d_T=w_*-w_T>0$ up to the earlier turn. In backward distance
$v=T_b-u$ put $Z(v)=\zeta(T_b-v)$. Then $Z(0)=0$ and for $v>0$,
\[
 (Z^2)_v=2\eps+2\phi'(T_b-v)Z .
\]
The right side extends continuously to $v=0$. Integrating first from
a positive lower endpoint and taking its limit gives this identity
in integral form from the turn. Put $c=1/[2(M+1)]$. While $Z\le c\eps$
the squared gap grows at speed at least $\eps$. It reaches $c^2\eps^2$
by $v=c^2\eps$ and cannot recross this level downwards. For $\eps$ so
small that $c^2\eps<b/2$, we have $\zeta(u)\ge c\eps$ from every
test turn back to the axis.

On the descent the same squared equation at $Z=\eps/\nu_*$ gives an
upper barrier starting from zero. Write $N(v)=\nu(T_b-v)$ there.
The error $e=Z-\eps/N$ obeys $e_v=-(N/Z)e-(\eps/N)_v$, with
$N/Z\ge\nu_*^2/\eps$ and $|(\eps/N)_v|\le M_1\eps/\nu_*^2$.
At $v=b/2$, $|e|\le C\eps$. Hence for $T_b-v\in I$,
$|e|\le C\eps e^{-\nu_*^2 b/(2\eps)}+C\eps^2=O(\eps^2)$.

The exact difference equation yields an integrable identity even at
the earlier turn,
\[
 (\log d_T)'=\frac{\eps}{\zeta(\zeta-d_T)}
             =\frac{\eps}{\zeta^2}+\frac{d_T'}{\zeta},\qquad d_T(T)=\zeta(T).
\]
Indeed $d_T'>0$, so its contribution at the turn satisfies
\[
 0\le E_T:=\int_0^T\frac{d_T'}\zeta\,\dd u
 \le\frac{\zeta(T)-d_T(0)}{\inf_{[0,T]}\zeta}=O(1).
\]
The numerator is at most $\zeta(T)=O(\eps)$ and the denominator is at
least $c\eps$. Integrating up to the turn by a one-sided limit gives
\[
 \log d_T(0)=\log \zeta(T)-\int_0^T\frac{\eps}{\zeta^2}\,\dd u-E_T .
\]
The same reference leg $w_*$ occurs for $T_0$, hence subtraction
cancels the integral on the common interval exactly. Also
$\mathfrak d(T)=d_T(0)[w_*(0)+w_T(0)]/2$. The upper-hit bounds and the bounds
on $\zeta(T)/\eps$ give bounded logarithms for both endpoint ratios. Thus
\[
 \log\frac{\mathfrak d(T)}{\mathfrak d(T_0)}
 =-\int_{T_0}^T\frac{\eps}{\zeta^2}\,\dd u+O(1).
\]
On $I$ the expansion for $\zeta$ gives
$\eps/\zeta^2=\nu^2/\eps+O(1)$. The integration length is at most
$L$, which gives the uniform remainder in \eqref{eq:deficit-ratio}.
\end{proof}


\begin{proposition}[Closing a finite incoming list with three additional powers]\label{prop:fold-closure}
Let $P_\lambda(x)=P_*(x)+\lambda x^{m_0}$, with fixed $P_*$ having zero coefficient at $m_0$, be a
monic polynomial of even degree $d_{\mathrm{top}}\ge 10$ with next term $b_1 x^{d_{\mathrm{top}}-1}$
($b_1>0$), an adjustable odd term $\lambda x^{m_0}$ ($3\le m_0\le d_{\mathrm{top}}-3$ odd), and all
remaining powers in $\{3,\dots,d_{\mathrm{top}}-2\}$. Suppose $\lambda\in I_\lambda$ varies on a
compact interval avoiding $0$ such that:
\begin{itemize}
  \item $\Delta(P_\lambda)=\ell_+(P_\lambda)-\ell_-(P_\lambda)$ has opposite strict signs at $\partial
  I_\lambda$ and a unique zero $\lambda_0\in\operatorname{int}(I_\lambda)$;
  \item $N_{\mathrm{ord}}$ pairwise disjoint compact ordinary sign brackets of $P_\lambda$ persist throughout $I_\lambda$;
  \item at $\lambda=\lambda_0$ and a fixed physical shift $s\in\mathbb{R}$, $h_1>h_2>\dots>h_M$ is a
  finite decreasing list of real numbers at which $G_+(h_r)-G_-(h_r+s)$ has nonzero signs
  $\sigma_r\in\{+1,-1\}$.
\end{itemize}
Then adjoining $\{d_{\mathrm{top}}+1,d_{\mathrm{top}}+3,d_{\mathrm{top}}+8\}$ produces, for all
sufficiently large outer dilations $R>1$, a polynomial $F_R(x)$ with
$\supp(F_R)=\supp(P_{\lambda_0})\cup\{d_{\mathrm{top}}+1,d_{\mathrm{top}}+3,d_{\mathrm{top}}+8\}$
whose half-return difference $\oscD_{F_R}(\beta)=b_+(\beta)-b_-(\beta)$ on $\{x=0,y>0\}$
simultaneously realizes, in strictly increasing order of $\beta$:
\begin{enumerate}
  \item the $N_{\mathrm{ord}}$ bounded ordinary sign brackets;
  \item an ascending fold sign bracket (signs $-,+$), a hyperbolic marked orbit on the fold ascent,
  and a descending fold sign bracket (signs $-,+$ ending at $\oscD_{F_R}(\beta_{\mathrm{refl}})>0$);
  \item points $\beta_{\mathrm{refl}}<\beta_1<\beta_2<\dots<\beta_M<\beta_{\mathrm{ext}}$ with
  $\sgn\oscD_{F_R}(\beta_r)=\sigma_r$ ($1\le r\le M$) and $\oscD_{F_R}(\beta_{\mathrm{ext}})>0$;
  \item a hyperbolic marked orbit whose launch is beyond $\beta_{\mathrm{ext}}$
  and which encloses all the orbits selected in the preceding items.
\end{enumerate}
After separating any shared strict endpoints, a small perturbation of the
$x^{m_0}$ coefficient gives a hyperbolic cycle in every sign bracket.
The two marked orbits continue within disjoint neighborhoods of their
original launches on the same support (Lemma~\ref{lem:hyperbolicity}).
\end{proposition}

The order of these launches is shown in Figure~\ref{fig:outer-fold}.

\begin{proof}
Shrink $I_\lambda$ around $\lambda_0$ if needed so the $M$ strict signs of
$G_{+,P_\lambda}(h_r)-G_{-,P_\lambda}(h_r+s)$ hold uniformly for $\lambda\in I_\lambda$. Put
$e=d_{\mathrm{top}}-2\ge 8$ and define on $r\ge 0$:
\begin{align*}
  E_{\mathrm{f}}(r) &= r^{d_{\mathrm{top}}} -
  \frac{d_{\mathrm{top}}}{d_{\mathrm{top}}+8}r^{d_{\mathrm{top}}+8}, & f_{\mathrm{f}}(r) &=
  \frac{E_{\mathrm{f}}'(r)}{r} = d_{\mathrm{top}} r^e(1-r^8), \\
  Q_{\mathrm{f}}(r) &= r^{d_{\mathrm{top}}-1}(1-r^2)(1-\rho r^2), & U_{\mathrm{f}}(a) &= \int_0^a
  (f_{\mathrm{f}}Q_{\mathrm{f}}' - f_{\mathrm{f}}'Q_{\mathrm{f}})\dd r, \quad \Tr_{\mathrm{f}}(t) =
  U_{\mathrm{f}}(t)-U_{\mathrm{f}}(1).
\end{align*}
\smallskip\noindent\emph{The two fold tests.} Let $z(y)>1$ denote the descending inverse of
$E_{\mathrm{f}}(z)=y<E_{\mathrm{f}}(1)=\frac{8}{d_{\mathrm{top}}+8}$, and set
$g=z(0)=(1+8/d_{\mathrm{top}})^{1/8}>1$. On $(1,\infty)$,
$Q_{\mathrm{f}}(r)/f_{\mathrm{f}}(r)=\frac{1}{d_{\mathrm{top}}}(A_0(r)-\rho B_0(r))$ where
$B_0(r)=(r^{-3}+r^{-1}+r+r^3)^{-1}$ satisfies $B_0'(r)<0$ (as $-3r^{-4}-r^{-2}+1+3r^2>0$ on
$(1,\infty)$) and $A_0(r)=B_0(r)/r^2$ satisfies $A_0'(r)<0$. Hence
$\rho\mapsto\Tr_{\mathrm{f}}(g)=\int_1^g f_{\mathrm{f}}(r)^2(Q_{\mathrm{f}}/f_{\mathrm{f}})'(r)\dd r$
is strictly increasing and affine with negative value at $\rho=0$, determining a unique $\rho>0$ with
$\Tr_{\mathrm{f}}(g)=0$. Writing $Q_{\mathrm{f}}(r)=(r^{e+1}-r^{e+3})-\rho(r^{e+3}-r^{e+5})$, direct
integration on $[0,1]$ gives
$\frac{1}{d_{\mathrm{top}}}\int_0^1[f_{\mathrm{f}}(r^p-r^{p+2})'-f_{\mathrm{f}}'(r^p-r^{p+2})]\dd r =
-4\bigl(\frac{e}{(e+p)(e+p+2)}-\frac{e+8}{(e+8+p)(e+10+p)}\bigr)$, whose sign is
$-\sgn(e(e+8)-p(p+2))$; this is negative at $p=e+1$ ($4e-3>0$) and positive at $p=e+3$ ($-15<0$), so
$U_{\mathrm{f}}(1)<0$ since $\rho>0$. Because $f_{\mathrm{f}}$ has negative leading coefficient at
$e+8$ and $Q_{\mathrm{f}}$ has positive leading coefficient $+\rho$ at $e+5$, $\Tr_{\mathrm{f}}(t)\sim
\frac{3d_{\mathrm{top}}\rho}{2e+13}t^{2e+13}\to+\infty$ as $t\to+\infty$.

On $(0,1)$, write $B_L(a)=U_{\mathrm{f}}(a)$ and
$B_R(a)=U_{\mathrm{f}}(a)-f_{\mathrm{f}}(a)\Tr_{\mathrm{f}}(z(E_{\mathrm{f}}(a)))/f_{\mathrm{f}}(z(E_{\mathrm{f}}(a)))$.
Since $\Tr_{\mathrm{f}}(g)=0$ and $z(y)$ is smooth near $y=0$,
$\Tr_{\mathrm{f}}(z(E_{\mathrm{f}}(a)))=O(a^{d_{\mathrm{top}}})$ as $a\downarrow 0$, whereas
$U_{\mathrm{f}}(a)/f_{\mathrm{f}}(a)\sim a^{d_{\mathrm{top}}-1}/(2e+1)>0$; as $a\uparrow 1$,
$U_{\mathrm{f}}(1)<0$ and $\Tr_{\mathrm{f}}(z)/f_{\mathrm{f}}(z)\to 0$ as $z\downarrow 1$. Fix
$0<a_0<a_1<a_*<1$ such that $(B_L,B_R)$ are both strictly positive at $a_0$ and both strictly negative
at $a_1$ and $a_*$, and fix $0<s_0<S_*$ large enough that $\Tr_{\mathrm{f}}(z(-s_0))>0$ and
$\Tr_{\mathrm{f}}(z(-S_*))>0$.

Writing $z_i=z(E_{\rm f}(a_i))$, the limiting order of the test radii is
\[
  a_0<a_1<a_*<1<z_1<z_0<g<z(-s_0)<z(-S_*).
\]
The descent near $g$ is thus available for turns at bounded physical
heights, between the reflected tests and the exterior one.

\smallskip\noindent\emph{Three finite matches.} In
$\operatorname{span}\{r^{d_{\mathrm{top}}+1},r^{d_{\mathrm{top}}+3}\}$, the linear map
$V\mapsto\bigl(V(1),\int_1^g(f_{\mathrm{f}}V'-f_{\mathrm{f}}'V)\dd r\bigr)$ is invertible because the
kernel of evaluation at $1$ is spanned by $r^{d_{\mathrm{top}}+1}(1-r^2)$, on which the second
functional is the negative of the nonzero $\rho$-slope of $\Tr_{\mathrm{f}}(g)$. Let
$V_1,V_2\in\operatorname{span}\{r^{d_{\mathrm{top}}+1},r^{d_{\mathrm{top}}+3}\}$ be the dual basis.
With $\xi=x/R$, define the physical primitive
\begin{align*}
  F_{R,\lambda,\theta_1,\theta_2}(x)
    ={}& P_\lambda(x) - \frac{d_{\mathrm{top}}}{d_{\mathrm{top}}+8} R^{-8}x^{d_{\mathrm{top}}+8} \\
       & + b_1 R^{d_{\mathrm{top}}-1}\bigl[Q_{\mathrm{f}}(\xi)-\xi^{d_{\mathrm{top}}-1}+\theta_1 V_1(\xi)+\theta_2 V_2(\xi)\bigr].
\end{align*}
Since $Q_{\mathrm{f}}(r)-r^{d_{\mathrm{top}}-1}=-(1+\rho)r^{d_{\mathrm{top}}+1}+\rho
r^{d_{\mathrm{top}}+3}$ has both coefficients strictly nonzero, for $|\theta_1|,|\theta_2|\le\delta$
sufficiently small all three outer powers
$\{d_{\mathrm{top}}+1,d_{\mathrm{top}}+3,d_{\mathrm{top}}+8\}$ retain nonzero coefficients. In scaled
coordinates $r=|x|/R$, $u=r^2/2$, $w=y/R^{d_{\mathrm{top}}}$, $\tilde{\tau}=R^{-1}$, and
$\tilde{\eps}=R^{2-2d_{\mathrm{top}}}$, the half-profiles take the form
$\phi_\sigma(u)=E_{\mathrm{f}}(r)+\sigma b_1\tilde{\tau}[Q_{\mathrm{f}}(r)+\theta_1 V_1(r)+\theta_2
V_2(r)]+o(\tilde{\tau})$.

For the estimates put $d=d_{\rm top}$ and fix $T_{\max}>z(-S_*)^2/2$.
On $\mathcal I=[0,T_{\max}]$ the omitted terms are bounded multiples of
$R^{m-d}(2u)^{m/2}$, $3\le m\le d-2$. Their sum is $O(R^{-2})$ in
$C^1(\mathcal I)$ uniformly on the parameter box. Each monomial slope is
monotone, so the slopes of the full profiles have uniformly bounded
variation on $\mathcal I$. On every $[\alpha,T_{\max}]$, $0<\alpha<T_{\max}$,
the profiles converge in $C^2$ to $E_{\rm f}(\sqrt{2u})$.
This preserves its simple crest and the signs of its slopes on compact
regular branches.

For each scaled lower axis depth $v\in[s_0,S_*]$, let $\widehat T_\sigma(v)$ be the unique turn
satisfying $Y_\sigma(\widehat T_\sigma(v))=v^2/2$ (uniqueness and existence follow from $Y_\sigma'>0$
and bracketing via \eqref{eq:envelope-bounds} on either side of $z(-v)$). Uniformly on $[s_0,S_*]$,
$\sqrt{2\widehat T_\sigma(v)}\to z(-v)$ and $\phi_\sigma(\widehat
T_\sigma(v))=-v+O(\sqrt{\tilde\eps})$, because the terminal ordinate is the suffix minimum back to the
axis. Let $\beta_a=\beta(Ra_*)$ be the upper launch of the plus turn
at physical radius $Ra_*$ on the ascent. At the common scaled lower
intercept $-S_*$ let $X_{\sigma,*}=X_\sigma(\widehat T_\sigma(S_*))$
be half the square of each scaled upper intercept. At bounded physical
height $h$ let $T_\sigma(h)$ near $\bar g=g^2/2$ be the descending
solution of $\phi_\sigma(T_\sigma(h))=h/R^d$. Its upper deficit from
that half's own deep reference is
$\mathfrak d_\sigma(h)=X_{\sigma,*}-X_\sigma(T_\sigma(h))>0$.
On $I_\lambda\times[-\delta,\delta]^2$ the equations and their
leading controls are
\begin{center}\small
\begin{tabular}{@{}ll ll@{}}
\toprule
Match & Control & Finite equality & Geometric effect\\
\midrule
(M1) & $\lambda$   & $b_+(\beta_a)=b_-(\beta_a)$
                 & Close the ascending orbit\\
(M2) & $\theta_1$ & $X_{+,*}=X_{-,*}$
                 & Close the deep reference orbit\\
(M3) & $\theta_2$ & $\log[\mathfrak d_+(0)/\mathfrak d_-(s)]=0$
                 & Anchor the shift on the descent\\
\bottomrule
\end{tabular}
\end{center}
Here $\lambda$ is the existing coefficient of $x^{m_0}$. The parameters
$\theta_1,\theta_2$ are independent combinations of the coefficients of
$x^{d+1},x^{d+3}$. The normalization of $V_1$ varies the crest value,
and that of $V_2$ varies the descending action at fixed level zero.
The table identifies the parameter controlling each pair of faces of
the box. We solve the three equalities simultaneously.

For (M1) choose fixed $0<t_1<a_*<t_0<1$. For each sign the upper
hits at turns $Rt_1,Ra_*,Rt_0$, divided by $R^{d_{\rm top}}$, tend
uniformly to the three strictly ordered positive values of $E_{\rm f}$
there, by \eqref{eq:envelope-bounds}. All these finite turns meet both
axes by Lemma~\ref{lem:tail-inward}(i). Monotonicity of $X_-$ thus
places the minus turn for the common launch between $Rt_1$ and $Rt_0$,
and the common physical launch tends to infinity.
For $0<r\le t_0$ the ratio $E_{\rm f}(r)/r^{d_{\rm top}}$ is at least
$1-d_{\rm top}t_0^8/(d_{\rm top}+8)>0$. On the physical interval
$K_0\le x\le Rt_0$ the odd terms and lower powers contribute
$O(x^{d_{\rm top}-1})$ uniformly. Hence
$F(\pm x)\ge c x^{d_{\rm top}}$ up to both turns.
Along any subsequence of parameters converging in the box to
$(\lambda',\theta_1',\theta_2')$ the folded primitives converge on
compacts in $C^1$ to $P_{\lambda'}$. The finite-turn assertion of
Lemma~\ref{lem:tail-inward}(ii) now gives both limiting intercepts
$\ell_\pm(P_{\lambda'})$. Continuity of the inner tail gap and compactness
give the uniform (M1) limit in \eqref{eq:miranda-box}.
For (M2) the upper hits at depth $S_*$ are within $O(\sqrt{\tilde\eps})$
of the perturbed maxima, uniformly on the same box. The maxima have
$r=1+O(\tilde\tau)$ and values
$E_{\mathrm f}(1)\pm b_1\tilde\tau\theta_1+O(\tilde\tau^2)$.
Dividing their upper-hit difference by $2b_1\tilde\tau$ gives
$\theta_1+o(1)$.

For (M3) the deficit is the integral of $X'_\sigma$ from
$T_\sigma(h)$ to that half's own reference turn $\widehat T_\sigma(S_*)$.
The order of these two turns makes the deficit positive throughout the
parameter box, including away from (M2). Since the reference tends to
$z(-S_*)^2/2<T_{\max}$, Lemma~\ref{lem:common-crest}\textup{(i)}
applies uniformly there and gives
\[
 -\tilde\eps\log \mathfrak d_\sigma(h)
 =\int_{r_{\rm crest,\sigma}}^{\sqrt{2T_\sigma(h)}}
       (\partial_r\phi_\sigma)^2\,\frac{\dd r}{r}
       +O(\sqrt{\tilde\eps}+\tilde\eps|\log\tilde\eps|),
\]
where $\partial_r$ differentiates $\phi_\sigma(r^2/2)$.
The descent through zero is regular. Changing the scaled terminal height
by $s/R^{d_{\rm top}}$ costs $O(R^{-d_{\rm top}})$ in the action.
At level zero its first variation in direction $V$ is
$\int_1^g(f_{\mathrm f}V'-f_{\mathrm f}'V)\dd r$.
It vanishes on $Q_{\mathrm f},V_1$ and equals one on $V_2$.
Together with the first two estimates this gives the following uniform
expansion on the box:
\begin{equation}\label{eq:miranda-box}
 \left(\oscD(\beta(Ra_*)), \frac{\sqrt{2X_{+,*}}-\sqrt{2X_{-,*}}}{2b_1\tilde\tau}, -\frac{\tilde\eps}{2b_1\tilde\tau}\log\frac{\mathfrak d_+(0)}{\mathfrak d_-(s)}\right)
 =\bigl(\Delta(P_\lambda),\theta_1,\theta_2\bigr)+o(1).
\end{equation}
Each continuous component therefore has opposite signs on its two faces
for all values of the other coordinates. Poincar\'e--Miranda gives
one common zero for every sufficiently large $R$. Every selection of
such zeros tends to $(\lambda_0,0,0)$ by compactness and uniqueness of
the limiting zero.

\smallskip\noindent\emph{The regular turns and the two matched orbits.} At
$(\lambda_R,\theta_{1,R},\theta_{2,R})$, the $N_{\mathrm{ord}}$ bounded ordinary brackets persist by
compact $C^1$ convergence to $P_{\lambda_0}$. Apply Lemma~\ref{lem:common-crest}\textup{(ii)} with
$E=E_{\mathrm f}$, $Q=Q_{\mathrm f}$, $\delta=b_1/R$ and the later
turns from (M2). Here $\theta_{i,R}\to0$ absorbs the two controls into
$o(\delta)$ in $C^1(\mathcal I)$, and $\sqrt{\tilde\eps}/\delta=R^{2-d}/b_1\to0$.
On the early ascent $K_0/R\le r\le t_0<1$, division of the radial
derivative by $r^{d-1}$ bounds each old lower power by a constant times
$(Rr)^{m-d}$. Choosing $K_0$ large preserves the rising slope there.
On $r\le K_0/R$ put $x=Rr$. The physical slope
$\partial_x F(\sigma x)/x$ is bounded since all powers are at least
three, so $\phi_\sigma'=O(R^{2-d})$ on an interval of $u$-length
$K_0^2/(2R^2)$. Its squared-slope integral is $O(R^{2-2d})=o(\delta)$.
Thus (M1) gives an ascending $-,+$ pair at $a_0,a_1$, the simple
orbit at $a_*$, and a descending $-,+$ pair at
$z(E_{\mathrm f}(a_1)),z(E_{\mathrm f}(a_0))$, ending at
$\beta_{\rm refl}$.

\smallskip\noindent\emph{The exterior test at a common lower hit.}
Set $v_0=s_0$, $V=S_*$, $\eps=\tilde\eps$ and $\alpha=b_1/R$.
In this calculation write $E,Q,f,\Tr$ for the outer fold data with
subscript $\mathrm f$. At the simultaneous match the profiles
$p_\pm(r)=\phi_\pm(r^2/2)$ have expansion $E\pm\alpha Q+o(\alpha)$
in $C^1$ from the crest through the depth-$S_*$ turns, as verified above.
Let $\widehat T_\pm(v)$ be the turns at the common scaled lower hit
$-v$, for $v_0\le v\le V$, and evaluate $J_\pm(v)$ on these actual
half-orbits. The radii $z(-v)$ lie on a compact regular descent.
For turns on a slightly larger strip the terminal level is negative
and bounded away from zero. The profiles converge uniformly to
$E\ge0$ before the crest and decrease on the outer descent, so the
terminal level is the minimum on the whole excursion. The lower envelope is constant.
By \eqref{eq:envelope-bounds} the lower hit differs from the terminal
level by $O(\sqrt\eps)$. Bracketing on this strip and $Y_\pm'>0$
give both turns, with terminal levels $-v+O(\sqrt\eps)$ uniformly.

The upper envelope is constant up to the perturbed crest and follows
the descent thereafter. In \eqref{eq:turn-derivs} the action errors
from replacing both legs by their envelopes are $O(\sqrt\eps)$,
by integration by parts and the uniform bounds on the slopes and
their variations. Replacing the terminal level by $-v$ in the descent
integral also costs $O(\sqrt\eps)$, since its level derivative is
$p_\pm'(r)/r$ there. Write $z_v=z(-v)$. At fixed level the endpoint
is $z_v\mp\alpha Q(z_v)/E'(z_v)+o(\alpha)$ and the moving crest
contributes zero to first order since $f(1)=0$. The first variation
is therefore $2\int_1^{z_v} fQ'\dd r-f(z_v)Q(z_v)=\Tr(z_v)$.
Regularity on the compact strip and $\sqrt\eps=o(\alpha)$ give, uniformly in $v$,
\[
 -\eps J_\pm(v)=C(v)\pm\alpha\Tr(z(-v))+o(\alpha),\qquad
 C(v)=\int_1^{z(-v)} f(r)^2r\dd r,\qquad C'(v)=-f(z(-v))>0 .
\]

Write $\beta_\pm(v)=R^d\sqrt{2X_\pm(\widehat T_\pm(v))}$ for the
physical upper hits and $K(v)=(\beta_+(v)^2-\beta_-(v)^2)/(2R^{2d})$.
The upper match gives $K(V)=0$. Differentiating
$Y_\pm(\widehat T_\pm(v))=v^2/2$ and using \eqref{eq:turn-derivs}
gives $K'(v)=v(e^{J_+(v)}-e^{J_-(v)})$ and $\beta_\pm'(v)>0$.
Let $\chi=\sgn\Tr(z(-v_0))$. Choose $v_0<v_1<v_2<V$ with
$\chi\Tr(z(-v))\ge c>0$ on $[v_0,v_2]$. There the smaller exponential
is at most $e^{-c\alpha/\eps}$ times the larger, with $\alpha/\eps\to\infty$.
Thus $-K'$ has sign
$\chi$ and $|K'|\ge(v_0/2)\max(e^{J_+},e^{J_-})$ for large $R$.
Put $\Delta=C(v_2)-C(v_1)>0$. Once
$|-\eps J_\pm-C|<\Delta/4$ on $[v_0,V]$, monotonicity of $C$ gives
\[
 \frac{\int_{v_2}^{V}|K'(v)|\dd v}{\int_{v_0}^{v_1}|K'(v)|\dd v}
 \le \frac{4V(V-v_2)}{v_0(v_1-v_0)}\,e^{-\Delta/(2\eps)} .
\]
Hence $K(v_0)=-\int_{v_0}^{V}K'(v)\dd v$ has sign $\chi$.
At $V$ the strict sign of $\Tr(z(-V))$ gives $K'(V)\ne0$ directly.

Both branches return up to the common upper hit
$\beta_*=\beta_+(V)=\beta_-(V)$, since every intervening finite turn
meets the axes and $X_\pm$ increase from zero. From
$b_\pm(\beta_\pm(v))=-R^d v$ we obtain
\[
 \oscD(\beta_+(v))=b_-(\beta_-(v))-b_-(\beta_+(v)).
\]
Since $b_-'<0$, this difference has the sign of $K(v)$.
At $\beta_*$ its derivative with respect to $v$ is
$\oscD'(\beta_*)\beta_+'(V)=-b_-'(\beta_*)[\beta_+'(V)-\beta_-'(V)]$.
The bracket equals $R^{2d}K'(V)/\beta_*\ne0$. The marked orbit is
therefore hyperbolic by Lemma~\ref{lem:variational}\textup{(i)--(ii)}.

Here $\Tr_{\mathrm f}(z(-s_0))>0$, so the exterior launch
$\beta_{\mathrm{ext}}=\beta_+(s_0)$ has $\oscD_{F_R}(\beta_{\mathrm{ext}})>0$.
The mark at $\beta_+(S_*)>\beta_{\mathrm{ext}}$ therefore encloses the
selected cycles at smaller launches.

\smallskip\noindent\emph{Bounded incoming heights.} We now compare turns near the physical radius
$Rg$. Lemma~\ref{lem:deficit-ratio} applies near $\bar g$ on a closed
strip of the strict descent terminating at the depth-$S_*$ turn.
Indeed $\widehat T_\sigma(S_*)\to z(-S_*)^2/2>\bar g$. We can take
a fixed compact test interval about $\bar g$ with a common buffer to
these reference turns, all below $T_{\max}$. The slope and its
derivative have the required bounds on the descent by $C^2$
convergence there. The upper hits tend to $E_{\mathrm f}(1)>0$.
Change variables from $u$ to $\kappa=\phi_\sigma(u)$ in
\eqref{eq:deficit-ratio}. For $h,h_0$ in any fixed compact set this gives
\begin{equation}\label{eq:physical-deficit}
 \tilde\eps R^{d_{\rm top}}\log\frac{\mathfrak d_\sigma(h)}{\mathfrak d_\sigma(h_0)}
 =-f_{\mathrm f}(g)(h-h_0)
     +O(R^{-1}+R^{2-d_{\rm top}}).
\end{equation}
Indeed $\kappa=h/R^{d_{\rm top}}$ at the upper limit,
$\phi_\sigma'(T_\sigma(h))=f_{\mathrm f}(g)+O(R^{-1})$, and
$\tilde\eps R^{d_{\rm top}}=R^{2-d_{\rm top}}$. The error here tends to zero on the physical height scale. By \textup{(M2)} the two marked upper hits agree, and by
\textup{(M3)} $\mathfrak d_+(0)=\mathfrak d_-(s)$. For every fixed $\eta>0$, the last
estimate brackets $\mathfrak d_+(h_r)$ between $\mathfrak d_-(h_r+s-\eta)$ and
$\mathfrak d_-(h_r+s+\eta)$. Write $B(h)=\beta(R\sqrt{2T_+(h)})$
for the actual launch height on the positive $y$-axis of the plus
trajectory that turns at physical ordinate $h$ on the far descent.
At $\beta_r=B(h_r)$ the minus turning ordinate is $h_r+s+o(1)$.
Write $a_R$ for either actual turning ordinate, with limit $a=h_r$
or $a=h_r+s$ on the corresponding half. All these ordinates belong
to a fixed compact height interval. Choose $1<\gamma_0<g$ close to
$g$. From the turn $x_T=R\sqrt{2T_\sigma(a_R)}$ down to $\gamma_0R$,
regular descent gives
$F_R(\sigma x)-a_R\ge cR^{d-1}(x_T-x)$. On the terminal layer
$x_T-x\le R^{-d}$ the graph equation gives
$0\le a_R-y_R(x)\le\sqrt{2x_T(x_T-x)}$. On the remaining part the
drop is bounded by integrating $CR^{2-d}/(x_T-x)$. Thus
$y_R(\gamma_0R)=a_R+O(R^{(1-d)/2}+R^{2-d}\log R)$ uniformly.

For $x/R\le\gamma_0$, $E_{\rm f}(x/R)/(x/R)^d\ge
1-(\gamma_0/g)^8>0$. The odd and lower powers contribute at most
$C x^{d-1}$ for $1\le x\le\gamma_0R$, uniformly in the matching box.
Hence for some $K_*$ and $c'>0$ we have $F_R(\sigma x)\ge c'x^d$
on $[K_*,\gamma_0R]$, and $P_{\lambda_0}(\sigma x)\ge c'x^d$ on
$[K_*,\infty)$. Let $H_*$ bound the absolute values of all turning ordinates under
consideration. Enlarge $K_*$ so that $c'x^d-H_*\ge c'x^d/2$
for $x\ge K_*$. Fix $K_1\ge K_*$. The further drop of $y_R$ down to $K_1$
and the tail loss of the limiting incoming graph $y_{\sigma,a}$ from
infinity to $K_1$ are each bounded by $C K_1^{2-d}$, with $C$
independent of $K_1,R$ whenever $\gamma_0R>K_1$.

On $[0,K_1]$ the continuous gap
$P_{\lambda_0}(\sigma x)-y_{\sigma,a}(x)$ has a positive minimum
over the two halves and the finite list of limiting heights.
Solve the folded graph equation backwards from the exact limiting
value $y_{\sigma,a}(K_1)$ and call the solution $Y_R$. For this fixed
$K_1$, compact convergence of $F_R$ and continuous dependence give
$Y_R\to y_{\sigma,a}$ uniformly, with $F_R(\sigma x)-Y_R(x)$ bounded
below by half the limiting minimum for large $R$. The two graphs $y_R,Y_R$ in the
same folded field contract their difference on backward integration.
Consequently
\[
 \limsup_{R\to\infty}|y_R(0)-y_{\sigma,a}(0)|\le 2C K_1^{2-d}.
\]
For any prescribed accuracy first choose $K_1$ to make this tail
bound small and then $R$ for that fixed interval. This proves
$b_+(\beta_r)\to G_{+,P_{\lambda_0}}(h_r)$ and
$b_-(\beta_r)\to G_{-,P_{\lambda_0}}(h_r+s)$, so the prescribed
strict signs persist.

Finally, write $x_{+,R}(h)=R\sqrt{2T_+(h)}$. For large $R$ the
actual plus turning radii satisfy
\[
 Rz_1<Rz_0<x_{+,R}(h_1)<\cdots<x_{+,R}(h_M)
 <R\sqrt{2\widehat T_+(s_0)}<R\sqrt{2\widehat T_+(S_*)}.
\]
At the ends this follows from $z_0<g<z(-s_0)<z(-S_*)$ and convergence
after division by $R$. All the middle radii divided by $R$ tend to
$g$. Their strict order follows from $h_1>\cdots>h_M$ and strict
descent of the actual plus profile. Strict increase of $X_+$ with
the turn gives the claimed order of upper launches.
\end{proof}

\begin{corollary}[Cost of transferring an ordered incoming list]\label{cor:sign-cost}
Under the hypotheses of Proposition~\ref{prop:fold-closure}, let
$\operatorname{var}(s_0,\ldots,s_n)$ count adjacent unequal signs.
For a finite incoming list with signs $\sigma_1,\ldots,\sigma_M$, the
same three new powers suffice for at least
\[
 N_{\rm ord}+4+\operatorname{var}(+,\sigma_1,\ldots,\sigma_M,+)
\]
simultaneous hyperbolic cycles on the enlarged support.
\end{corollary}
\begin{proof}
The two regular fold brackets and the two simple marked orbits are
disjoint from the incoming brackets. The positive return at each end
of the incoming list accounts for the two exterior signs in the formula.
Shrink adjacent brackets at a shared strict sign and apply
Lemma~\ref{lem:hyperbolicity} to the existing coefficient of $x^{m_0}$.
\end{proof}
